\clearpage
\fancyhead[L]{\small PALLAS. CERES. VESTA.}
\begin{center}
\textbf{PALLAS. CERES. VESTA.}
\par\vspace{4pt}\hrule\vspace{4pt}
\fontsize{10}{12}\selectfont Göttingische gelehrte Anzeigen. Stück 136. Seite 1348..1350. 1809 August 26.
\par\vspace{4pt}\hrule\vspace{4pt}
\end{center}

\par\indent Eine kurze Nachricht in unseren Blättern über die Wiederauffindung der \emph{Pallas} an der hiesigen Sternwarte in der Nacht vom 9. auf den 10. August dürfte den Astronomen willkommen sein. Der Mangel an brauchbaren Beobachtungen aus dem vergangenen Jahr hatte Herrn Professor \textsc{Gauss} davon abgehalten, die Elemente zu verbessern. Bei der Suche nach dem Planeten muss man sich daher noch an die älteren Elemente halten, nach denen Herr Professor \textsc{Bode} im astronomischen Jahrbuch für 1811 eine kleine Ephemeride berechnet hat. Schon am 8. August wurde die Gegend abgesucht, in der Pallas zu erwarten war. Am folgenden Abend hatte sich einer der aufgezeichneten Sterne beträchtlich nach Süden bewegt. Am 10. August bestätigte es sich vollständig, dass dies Pallas gewesen war. Herr Professor \textsc{Gauss} hat bisher die folgenden zwei Beobachtungen gemacht [dieselben wie die an von \textsc{Lindenau} im Brief vom 14. Aug. 1809 mitgeteilten, S. 314 dieses Bandes]. Die erste Deklination ist nicht ganz so zuverlässig wie die zweite. Beide Rektaszensionen können als gut gelten. Demnach gibt die erwähnte Ephemeride von Herrn Professor \textsc{Bode} die Deklination gut an, die Rektaszension aber um 20 Minuten zu klein. Pallas erreicht übrigens schon die volle Helligkeit eines Sterns neunter Größe.

\par\indent Am 10. August wurde auch \emph{Ceres} wiedergefunden. Herr Professor \textsc{Gauss} hat sie einige Male mit benachbarten Sternen verglichen. Diese Beobachtungen sind aber noch nicht vollständig reduziert. Vorerst genügt es festzustellen, dass der Planet auf die Minute genau an der Position gefunden wurde, an der man ihn nach den XIII. Elementen (siehe das 87. Stück der Gött. gel. Anz. dieses Jahres) [3. Juni 1809, S. 313 dieses Bandes] und der danach berechneten Ephemeride (im Maiheft der Monatl. Correspondenz) erwartet hatte. Ceres erscheint ebenfalls so hell wie ein Stern neunter Größe.

\par\indent Auch \emph{Vesta} glaubte Herr Professor \textsc{Harding} bereits am 12. August in einem schwachen Sternchen neunter Größe zu erkennen. Es stand sehr nahe bei der Position, an der sie nach den IV. Elementen (siehe das 55. Stück dieser Anz.) [8. April 1809, S. 311 dieses Bandes] und der danach von Herrn Dr. \textsc{Schumacher} berechneten Ephemeride (Monatl. Corresp. am angegebenen Ort) stehen sollte. Die Witterung hat seither allerdings noch keine genauere Bestätigung zugelassen.
